\documentclass[12pt]{article}
\usepackage[margin=1in]{geometry}
\usepackage{enumitem}
\usepackage{titlesec}
\usepackage{parskip}
\usepackage{amsmath}
\usepackage{hyperref}
\usepackage{fontenc}

\title{\textbf{North 1990 Claude Guide}\\
\large North, Douglass C. \textit{Institutions, Institutional Change and Economic Performance}.\\
Cambridge: Cambridge University Press, 1990.}
\author{}
\date{}

\begin{document}

\maketitle
\tableofcontents
\newpage

%--------------------------------------------------
\section{Central Claims}
%--------------------------------------------------

\begin{itemize}[leftmargin=*, itemsep=4pt]
    \item \textbf{Institutions defined}: ``the humanly devised constraints that structure political, economic, and social interaction.'' They consist of \textbf{formal rules} (constitutions, laws, property rights), \textbf{informal constraints} (norms, conventions, codes of conduct), and their \textbf{enforcement characteristics}. Institutions reduce uncertainty by providing a stable (if not necessarily efficient) structure to everyday economic and political life.
    \item \textbf{Institutions vs.\ organizations}: institutions are the ``rules of the game''; organizations (firms, political parties, unions, agencies) are the ``players.'' Organizations emerge and evolve in response to the incentive structure institutions create, and in turn become the vehicles through which institutional change occurs.
    \item \textbf{Central puzzle}: neoclassical economics assumes frictionless exchange and predicts convergence toward efficient outcomes, yet real economies persistently exhibit poor-performing institutions and fail to converge on efficient arrangements. North's explanation: exchange is costly (\textbf{transaction costs}---the costs of measuring what is being exchanged and of enforcing agreements), and institutions exist precisely to economize on these costs. Where institutions fail to reduce transaction costs effectively, economic performance suffers, often persistently.
    \item \textbf{Path dependence}: institutional change is path-dependent. Increasing-returns mechanisms (analogous to Brian Arthur's and Paul David's work on technological lock-in, e.g.\ QWERTY) mean that early institutional choices constrain and channel later development, even toward inefficient long-run trajectories, because switching costs rise as organizations and complementary institutions adapt to and reinforce the status quo framework.
    \item \textbf{Adaptive efficiency vs.\ allocative efficiency}: North argues that what matters most for long-run growth is not static allocative efficiency (the neoclassical focus on optimal resource allocation given fixed institutions) but \textbf{adaptive efficiency}---an institutional framework's capacity to encourage decentralized experimentation, tolerate and learn from failure, and adjust over time to changing circumstances.
    \item \textbf{Ideology and mental models}: institutional persistence and change depend not only on objective incentive structures but on the shared \textbf{mental models}/ideologies actors use to interpret the world, which shape perceptions of fairness, legitimacy, and compliance---and which themselves change only slowly, contributing to institutional path dependence even after formal rules are altered.
\end{itemize}

%--------------------------------------------------
\section{Transaction Costs and the Theory of Institutions}
%--------------------------------------------------

\begin{itemize}[leftmargin=*, itemsep=4pt]
    \item \textbf{Two components of transaction costs}: (1) \textit{measurement costs}---the cost of ascertaining the valuable attributes of the good, service, or performance being exchanged; (2) \textit{enforcement costs}---the cost of protecting rights and enforcing agreements against opportunism, cheating, or default.
    \item \textbf{Personal vs.\ impersonal exchange}: in small, face-to-face communities, informal norms, reputation, and repeated interaction can enforce agreements cheaply. As economies scale up toward large, impersonal, long-distance exchange, informal enforcement becomes insufficient, and formal third-party enforcement (courts, contract law, state-backed property rights) becomes necessary---but is itself costly to build and sustain.
    \item \textbf{Property rights}: secure, well-specified, and reliably enforced property rights are the central institutional precondition for productive economic exchange and investment; insecure property rights raise the effective cost of economic activity and depress investment, innovation, and growth.
    \item \textbf{The Coasean foundation}: North's transaction-cost framework builds directly on Ronald Coase's insight that exchange is costly and that the allocation of legal rights (property rights, liability rules) therefore has real economic consequences---in a zero-transaction-cost world (the Coase theorem's counterfactual), institutions would not matter; because transaction costs are always positive, institutions matter enormously.
    \item \textbf{Historical illustrations}: North draws on economic history---the medieval Law Merchant and merchant guilds that emerged to enforce long-distance trade agreements privately before strong states existed, and the comparative development of property-rights institutions in early modern Western Europe versus other regions---to show institutions evolving specifically to reduce the transaction costs of increasingly complex and impersonal exchange.
\end{itemize}

%--------------------------------------------------
\section{The State, Credible Commitment, and Institutional Change}
%--------------------------------------------------

\begin{itemize}[leftmargin=*, itemsep=4pt]
    \item \textbf{The paradox of the state}: the state is essential for enforcing property rights and contracts at a scale informal mechanisms cannot reach, yet the same coercive power that lets the state enforce rights also lets it violate them---expropriating wealth, reneging on commitments, or granting predatory monopolies to favored actors. Economic growth therefore depends on solving a \textbf{credible commitment problem}: how can a state powerful enough to enforce property rights be constrained from abusing that same power?
    \item \textbf{Constitutional and institutional constraints on the state}: North's answer (elaborated further in North and Weingast's related work on England's Glorious Revolution of 1688) is that durable economic growth requires institutions---representative assemblies, independent judiciaries, constitutional limits---that credibly constrain the sovereign's ability to arbitrarily alter property rights or default on debts.
    \item \textbf{Incremental vs.\ discontinuous change}: formal rules (laws, constitutions) can change discontinuously, often through political action or crisis, but informal constraints (norms, conventions, embedded mental models) change only gradually. This mismatch explains why formal institutional reforms (e.g., transplanting Western legal codes or constitutions into other contexts) frequently fail to produce the expected economic performance---the informal ``software'' has not caught up with the formal ``hardware.''
    \item \textbf{Institutional entrepreneurs}: change is driven by political and economic entrepreneurs who perceive that altering the rules would be profitable given shifts in relative prices, technology, or bargaining power, but they operate within, and are constrained by, the existing institutional matrix and the organizations it has produced.
    \item \textbf{Organizations reinforce institutions}: because organizations invest in skills and strategies suited to the existing institutional framework (whether that framework is efficient or not), they develop a vested interest in its perpetuation---a key mechanism generating path dependence and explaining the persistence of inefficient institutional arrangements over long historical periods.
\end{itemize}

%--------------------------------------------------
\section{Core Works Cited}
%--------------------------------------------------

\begin{itemize}[leftmargin=*, itemsep=4pt]
    \item \textbf{Ronald Coase}, ``The Nature of the Firm'' (1937) and ``The Problem of Social Cost'' (1960)---the foundational transaction-cost economics that underlies North's entire framework.
    \item \textbf{Oliver Williamson}, \textit{The Economic Institutions of Capitalism} (1985)---transaction-cost analysis of firms, contracts, and opportunism; a close contemporary complement to North's institutional economics.
    \item \textbf{Armen Alchian and Harold Demsetz}, ``Production, Information Costs, and Economic Organization'' (1972)---property rights and the theory of the firm.
    \item \textbf{Paul David}, ``Clio and the Economics of QWERTY'' (1985), and \textbf{W.\ Brian Arthur}'s work on increasing returns---the technological path-dependence literature North adapts to institutional change.
    \item \textbf{Douglass North and Robert Thomas}, \textit{The Rise of the Western World} (1973)---North's earlier collaborative economic history, laying groundwork for the institutional-transaction-cost account of European growth developed more fully here.
    \item \textbf{Douglass North and Barry Weingast}, ``Constitutions and Commitment'' (1989)---the credible-commitment account of England's Glorious Revolution, directly extending this book's state/commitment framework.
\end{itemize}

%--------------------------------------------------
\section{Methodological Contributions and Limitations}
%--------------------------------------------------

\begin{itemize}[leftmargin=*, itemsep=4pt]
    \item \textbf{New Institutional Economics}: the book is a foundational text of the ``New Institutional Economics'' school, bridging economics and comparative politics by making institutions (rather than technology or factor endowments alone) central to explaining cross-national and cross-historical variation in economic performance.
    \item \textbf{Bridging economics and history}: North's method combines neoclassical price-theoretic reasoning with deep historical and comparative case evidence, an approach that influenced both economic history and comparative political economy.
    \item \textbf{Critiques}: critics have argued the framework is difficult to falsify (institutions can be invoked post hoc to explain almost any growth outcome), that ``path dependence'' risks becoming a residual category explaining persistence without precise predictive content, and that the framework underspecifies exactly how and when informal norms change or catch up to formal rules.
    \item \textbf{Legacy}: the book's core ideas (institutions as rules of the game, transaction costs, path dependence, credible commitment) became central to subsequent work across economics, comparative politics, and development studies, including the World Bank's ``good governance'' agenda and the broader institutionalist turn in comparative political economy.
\end{itemize}

%--------------------------------------------------
\section{Connections to the Broader Literature}
%--------------------------------------------------

\begin{itemize}[leftmargin=*, itemsep=4pt]
    \item \textbf{Ostrom 1990}: both books treat institutions as solutions to underlying cooperation and transaction-cost problems, but at different scales---Ostrom focuses on community-level, self-governed common-pool-resource institutions, while North focuses on macro-level, often state-backed institutions governing impersonal market exchange and property rights.
    \item \textbf{Olson 1965}: North's transaction-cost logic and Olson's collective-action logic are complementary---both explain why voluntary, decentralized action alone often fails to produce efficient outcomes, requiring institutional solutions (whether selective incentives or property-rights enforcement) to align individual and collective interest.
    \item \textbf{Huntington 1968}: Huntington's concept of political institutionalization (the acquisition of value, stability, and complexity by political organizations and procedures) parallels North's account of institutions providing stability and structure, though Huntington emphasizes managing political participation while North emphasizes reducing economic transaction costs.
    \item \textbf{Weber 1968 / Weber 2013}: North's formal, rule-bound institutions enforced by predictable third parties (courts, states) are the economic-historical analogue of Weber's legal-rational authority; both frameworks treat predictable, depersonalized rule-following as the foundation of modern large-scale coordination, in contrast to more personalistic/traditional arrangements.
    \item \textbf{Cox 1987 / Cox 1997}: Cox's account of British parliamentary and party institutions evolving in response to changing electoral incentives is a case-specific instance of the general North-style logic of institutional entrepreneurs adapting rules to changing relative costs and benefits.
    \item \textbf{Putnam 1993}: North's institutions (formal and informal rules enforced by the state or by convention) and Putnam's social capital (informal trust and networks) are two distinct but related solutions to the same underlying problem of enabling cooperation and exchange beyond narrow kin or patron-client ties; Putnam's account can be read as elaborating the ``informal constraints'' side of North's institutional framework in much greater sociological depth.
    \item \textbf{Przeworski et al.\ 2000}: North's theoretical claims about institutions and economic performance are the direct conceptual backdrop against which Przeworski et al.'s empirical findings on regime type, institutions, and growth can be read---testing, in a specific domain (democracy vs.\ dictatorship), the broader institutional-performance relationship North theorizes.
\end{itemize}

\end{document}
